% !PRINTED 744
\weberpage{744}
%<<CONT
nach \S.\,178, 1. das Product einer reellen Einheit und einer
Einheitswurzel $r^{2\lambda}$ sein, und da auch $\Delta(r)$
reell ist, so müsste $r^{2\lambda+1}$ reell sein, was nicht
möglich ist. Demnach ist $\Delta(r)$ in $H_\mu$ enthalten.

Da nun $\mu$ ebenso wie $m$ jede beliebige Potenz von 2 sein
kann, so ergiebt sich hieraus das zweite Haupttheorem:

B. \gesperrt{Eine Einheit des Körpers $H_m$, die mit allen
ihren Conjugirten positiv ist, ist das Quadrat einer Einheit
desselben Körpers.}

Von den beiden Theoremen A. und B. haben wir aber im \S.\,184
den Beweis des Hauptsatzes von den \gesperrt{Abel}'schen Zahlkörpern
(\S.\,179, I.) abhängig gemacht.

\begin{center}
\rule{4cm}{0.4pt}
\end{center}
